\begin{abstract}

Nesse trabalho buscamos determinar o valor da carga fundamental $e$ 
utilizando uma metodologia similar à utilizado por R. Millikan e H. Fletcher
em 1909.
O experimento consiste em borrifar gotícular de óleo no espaço entre dois capacitores
de placas paralelas. 
As gotas tem chance de sofrer eletrização por atrito ao serem borifadas, e assim, 
respondem ao campo elétrico entre as placas do capacitor, que tem valor conhecido. 
Assim, medindo a velocidade com que as gotas sobem e descem, foi possível 
determinar a carga e o raio de cada uma delas. 
Com os dados obtidos, é possível observar que as cargar tendem a ocorrem em 
múltiplos inteiros de um valor fundamental, que foi calculado. 

\end{abstract}
